\section{Ablations}\label{sec-8}

\textbf{Pre-registered ablations}:

\begin{table*}[t]
\centering\scriptsize\setlength{\tabcolsep}{3.5pt}
\caption{Pre-registered ablations A1--A12 and their status.}\label{tbl:ablations}
\begin{tabular}{@{}
>{\raggedright\arraybackslash}p{(\linewidth - 4\tabcolsep) * \real{0.3333}}
  >{\raggedright\arraybackslash}p{(\linewidth - 4\tabcolsep) * \real{0.3333}}
  >{\raggedright\arraybackslash}p{(\linewidth - 4\tabcolsep) * \real{0.3333}}@{}}
\toprule
\begin{minipage}[b]{\linewidth}\raggedright
\#
\end{minipage} & \begin{minipage}[b]{\linewidth}\raggedright
Name
\end{minipage} & \begin{minipage}[b]{\linewidth}\raggedright
Status
\end{minipage} \\
\midrule
A1 & Top-1 vs Top-5 metric A/B & partial (Phase 4a tables include both) \\
A2 & Strict ID vs cross-mapped ID & superseded by A4 \\
A3 & Backbone \ensuremath{\times} Scaffolding 2\ensuremath{\times}N & done (\Cref{sec-6-2} / Phase 4a matrix) \\
A4 & Strict vs ORPHA-fuzzy-variants & \textbf{done --- this section \Cref{sec-8-1}} \\
A5 & Silver gold vs physician gold & \textbf{interim done} (Opus 4.8 agent gold; \Cref{sec-9} L5) --- physician swap at camera-ready \\
A6 & TS-Guessing contamination audit & \textbf{done --- \Cref{sec-8-9} + \Cref{sec-7-10}} \\
A7 & Single LLM judge vs dual-judge (Gemini\ensuremath{\rightarrow}Claude) & partial (\Cref{sec-7-5} done) \\
A8 & Reasoning on vs off (thinking-mode, = H6) & \textbf{done --- \Cref{sec-8-10}} (V4-Pro on/off) \\
A9 & Subprocess timeout cap 300s vs 600s vs 1200s & \textbf{done --- \Cref{sec-8-6}} \\
A10 & Prevalence-stratified R@1 & done (\texttt{A10\_prevalence\_stratified.md}) \\
A11 & Cross-dataset agent ranking stability & done (\Cref{sec-7-6} / \texttt{A11\_ranking\_stability.md}) \\
A12 & LLM-judge swap (P5 with Claude vs Gemini judge) & done (\Cref{sec-7-5}) \\
\bottomrule
\end{tabular}
\end{table*}

\subsection{Ablation A4 --- Strict vs ORPHA-Fuzzy Variants Cross-Map}\label{sec-8-1}

\textbf{Question}: When an LLM emits a generic disease name (e.g.~``Methylmalonic
Acidemia'') that fuzzy-matches multiple ORPHA codes at tied scores
(ORPHA:26 ``Methylmalonic acidemia with homocystinuria'', ORPHA:27
``Vitamin B12-unresponsive methylmalonic acidemia'', ORPHA:280183 ``\ldots{}
transcobalamin receptor defect''), does the evaluator credit the
prediction?

\textbf{A. Strict baseline}: prediction must exact-match a gold OMIM / ORPHA /
CCRD ID or cross-map to one via Orphadata.

\textbf{B. Variants-aware}: adapter logs the \emph{tied top-K} (score \ensuremath{\geq} top - 5)
ORPHA candidates in \texttt{extra{[}"ranked\_predictions\_variants"{]}}; evaluator
returns True if any tied variant hits gold.

\textbf{Legacy result audit} (Phase 4a originally ran N=100 \ensuremath{\times} 4 datasets \ensuremath{\times} 5
backbones; the table below reports only the three diagnostic datasets):

\begin{table*}[t]
\centering\scriptsize\setlength{\tabcolsep}{3.5pt}
\caption{ORPHA-variant evaluation-channel effect (A9).}\label{tbl:a9}
\begin{tabular}{@{}
lll@{}}
\toprule
Dataset & Aggregate \ensuremath{\Delta} R@1 & Aggregate \ensuremath{\Delta} R@5 \\
\midrule
Phenopacket-Store & \textbf{+0.03} & \textbf{+0.06} \\
RareArena RDS & \textbf{+0.02} & \textbf{+0.04} \\
RareBench HF & +0.00 & +0.00 \\
\textbf{All diagnostic datasets combined} & \textbf{+0.010} & \textbf{+0.020} \\
\bottomrule
\end{tabular}
\end{table*}

\textbf{Top-impacted cells}:
- PP-Store mdagents: +3 pp R@1, +6 pp R@5
- PP-Store medagents: +4 pp R@1, +6 pp R@5
- PP-Store llm\_control: +3 pp R@1, +6 pp R@5

\textbf{Where it doesn't help}:
- vc\_rdagent / lirical: bypass name-mapping (use IDs directly)
- RareBench: gap is ORPHA hierarchy-level, beyond fuzzy-tie scope
- DeepRare: emits domain-specific name spellings that fuzzy already handles

\textbf{Decision}: variants-aware metric is on by default in main \Cref{tbl:main};
strict variant reported in parentheses for reviewer reference.

\subsection{Ablation A3 --- Backbone \ensuremath{\times} Scaffolding Cross-Product}\label{sec-8-2}

Detailed in \Cref{sec-6-2} main results. Highlights:

\begin{itemize}
\tightlist
\item
  \textbf{No-scaffold control} (llm\_control) is backbone-insensitive on
  PP-Store (R@1 = 0.27-0.30 across 4 backbones, full-N; V4-Pro reasoning-off)
\item
  \textbf{Single-pass multi-agent} (mdagents) has a narrow 4-pp backbone spread
  (Gemini 0.28, DS V4-Pro-off 0.27, V4-Flash 0.25, GPT-5 0.24)
\item
  \textbf{Multi-round debate} (medagents) tops on Gemini (0.30), then V4-Pro-off /
  GPT-5 (0.28), weakest on V4-Flash (0.25)
\item
  \textbf{OSCE simulation} (agentclinic) collapses on GPT-5 minimal (0.13)
\item
  \textbf{Panel orchestration} (MAI-DxO) collapses universally on HPO input
\end{itemize}

The cross-product matrix shows no universally-best backbone: each
scaffolding architecture has its own preferred backbone. This refutes
the ``just use the best backbone'' simplifying assumption.

\subsection{Ablation A9 --- Subprocess Timeout Cap (300s vs 600s vs 1200s)}\label{sec-8-6}

The subprocess wall-clock cap interacts with two distinct failure modes,
which A9 disentangles:

\textbf{(a) Borderline-slow but legitimate cells --- cap matters.} During the
N=500 rerun, \texttt{medagents\ \ensuremath{\times}\ DS\ V4-Flash} and \texttt{agentclinic\ \ensuremath{\times}\ DS\ V4-Flash}
on RareBench and the legacy MIMIC ICD-title task showed many \texttt{timeout} records at the default 300s
cap. A probe re-ran a representative medagents timeout case at a 900s
cap and it completed successfully in 309s --- i.e.~the case was not
hung, the 300s cap was simply slightly too tight (compounded by the
empty-content retry adding extra subprocess invocations). Raising the cap to 600s and
re-running recovered these cells essentially completely (agentclinic
RareBench 60/60 ok; legacy MIMIC agentclinic 37/37 and medagents 267/268). DS
V4-Pro remained slow enough that 600s still left a small timeout tail
(mdagents V4-Pro RareBench 21/36 recovered) --- genuine backbone latency,
reported honestly.

\textbf{(b) Genuinely degenerate output --- cap does NOT help.} MAI-DxO \ensuremath{\times} GPT-5
still degenerates at the 1200s cap (\Cref{sec-9} L1): the panel emits no usable
ranked diagnosis regardless of time budget. This is an architecture\ensuremath{\times}backbone
incompatibility, not a latency problem; more wall-clock buys nothing.

\textbf{Conclusion}: the timeout cap is a real evaluation hyperparameter for
slow-but-valid cells (600s is the right default for hosted DeepSeek/Gemini
on long free-text), but it cannot rescue genuinely degenerate
agent\ensuremath{\times}backbone pairings. Cap choices are logged per-cell.

\subsection{Ablations deferred / data-gated}\label{sec-8-7}

\begin{itemize}
\tightlist
\item
  \textbf{A5} (silver gold vs physician gold) --- pending 200-case PMC OA
  holdout physician annotation (handoff package prepared; annotation in
  progress).
\item
  \textbf{A6} (TS-Guessing contamination audit) --- n-gram overlap vs LLM
  training cutoff; gated on the post-cutoff holdout being curated.
\item
  \textbf{A10} (prevalence-stratified R@1) --- computed.
\item
  \textbf{A11} (cross-dataset ranking stability) --- computed (see \Cref{sec-7-6}).
\end{itemize}

\subsection{Pre-registration check}\label{sec-8-5}

Under our research-integrity protocol, all hypotheses (H1-H11) and
ablations (A1-A12) were frozen at OSF prior to Phase 4a launch. The
above results are the first pre-registered evaluation report.
\textbf{No post-hoc hypothesis selection.}

\subsection{Holm-Bonferroni family-wise correction over H1-H11}\label{sec-8-8}

Per pre-registration we apply Holm-Bonferroni at \ensuremath{\alpha}=0.05 (one-sided in the
predicted direction) over the testable subset of H1-H11. Family size m=6
(H3/H5/H9 excluded --- data unavailable, see \Cref{sec-9} L3/L4 and tasks \#63/\#64/\#66; H6
now tested descriptively in \Cref{sec-8-10} but kept out of this z/\ensuremath{\rho} family). 2026-07-06 full-N refresh: all inputs recomputed on full-N; H2 upgraded from an n=50
pilot to an n=500 paired design.

\begin{table*}[t]
\centering\scriptsize\setlength{\tabcolsep}{3.5pt}
\caption{Holm--Bonferroni family-wise correction over H1--H11.}\label{tbl:holm}
\begin{tabular}{@{}
>{\raggedright\arraybackslash}p{(\linewidth - 10\tabcolsep) * \real{0.1667}}
  >{\raggedright\arraybackslash}p{(\linewidth - 10\tabcolsep) * \real{0.1667}}
  >{\raggedright\arraybackslash}p{(\linewidth - 10\tabcolsep) * \real{0.1667}}
  >{\raggedright\arraybackslash}p{(\linewidth - 10\tabcolsep) * \real{0.1667}}
  >{\raggedright\arraybackslash}p{(\linewidth - 10\tabcolsep) * \real{0.1667}}
  >{\raggedright\arraybackslash}p{(\linewidth - 10\tabcolsep) * \real{0.1667}}@{}}
\toprule
\begin{minipage}[b]{\linewidth}\raggedright
\#
\end{minipage} & \begin{minipage}[b]{\linewidth}\raggedright
Claim
\end{minipage} & \begin{minipage}[b]{\linewidth}\raggedright
Stat
\end{minipage} & \begin{minipage}[b]{\linewidth}\raggedright
raw p
\end{minipage} & \begin{minipage}[b]{\linewidth}\raggedright
Holm-adj p
\end{minipage} & \begin{minipage}[b]{\linewidth}\raggedright
Survives \ensuremath{\alpha}=0.05?
\end{minipage} \\
\midrule
\textbf{H1} & Classical \textgreater{} LLM R@1 on super-rare tier (\textless1/1,000,000) & z=17.54 & 3.7e-69 & 2.2e-68 & \textbf{yes} \\
\textbf{H8} & R@1 at 16-30 HPO terms \textgreater{} \ensuremath{\leq}5 (inverted-U left tail) & z=12.57 & 1.6e-36 & 7.8e-36 & \textbf{yes} \\
\textbf{H2} & llm\_control P3 \textgreater{} P2 (genotype channel, full-N paired) & z=6.40 & 7.6e-11 & 3.0e-10 & \textbf{yes} \\
\textbf{H7} & Cross-agent specialty rank \ensuremath{\rho} \textgreater{} 0.6 & \ensuremath{\rho}=0.92 & 5.5e-04 & 1.6e-03 & \textbf{yes} \\
\textbf{H4} & Scaffold benefit larger on multi-system than single (DoD) & z=2.61 & 4.5e-03 & 9.0e-03 & \textbf{yes} \\
H10 & Spearman \ensuremath{\rho}(faithfulness, accuracy) \textless{} 0.5 & \ensuremath{\rho}=0.35 & 3.7e-02 & 3.7e-02 & nominal but judge-dependent (see below) \\
\bottomrule
\end{tabular}
\end{table*}

\textbf{Reading}: 5 of 6 testable hypotheses now survive the strictest
pre-registered family-wise correction (up from 2/6 at pilot N). The full-N
reruns flipped H2, H4, and H7 from under-powered to significant: the
genotype-channel lift (H2: +19.8 pp, n=500 paired, McNemar \ensuremath{\chi}\textsuperscript{2}=85), the
scaffold-benefit-on-complexity difference-of-differences (H4), and the
cross-agent specialty blind-spot correlation (H7: \ensuremath{\rho}=0.92 across 18 specialties)
are all now FWE-robust, alongside the two headline claims H1 (classical dominate
super-rare) and H8 (interior phenotype-density optimum). H10
(faithfulness--accuracy decoupling) we report as exploratory, not a clean rejection. An earlier version rested on a family-judge (Gemini) \ensuremath{\rho}=0.098 vs
non-family (Claude) \ensuremath{\rho}=0.616 split, but a 2026-07-22 frozen-audit re-run showed
that split was a trace-capture artifact: the Gemini scores had been computed
on truncated/empty traces while the Claude scores used repaired traces. Re-running
the Gemini judge on the \emph{same} repaired traces (n=40) raises its \ensuremath{\rho} from 0.098 to
\textbf{0.457}, against Claude's 0.640 (judge agreement \ensuremath{\rho}=0.741) --- a modest residual
judge-family difference, not the near-zero ``strong decoupling'' originally claimed.
The pre-registered H10 verdict (\ensuremath{\rho}\textless0.5) is therefore genuinely borderline and
judge-sensitive, so we keep H10 exploratory and withdraw the \ensuremath{\rho}=0.098 figure. See
\Cref{sec-7-5} for the de-confounded analysis.

\subsection{Ablation A6 --- TS-Guessing data-contamination audit}\label{sec-8-9}

We probe anticipated objection \#1 (``LLMs answer well only on diseases they were
trained on'') by correlating, for each gold ORPHA in our phase4a predictions,
the pre-cutoff PubMed mention count with per-disease R@1, separately per
backbone. PubMed query uses \texttt{esearch.fcgi} with \texttt{"\textless{}disease\ name\textgreater{}"{[}All\ Fields{]}}
and \texttt{maxdate=2024/06/30} (conservative cutoff covering all four backbones).
Spearman \ensuremath{\rho} over log(mention + 1) vs R@1; per-disease aggregate requires
\ensuremath{\geq}3 predictions per (disease, backbone).

\begin{table*}[t]
\centering\scriptsize\setlength{\tabcolsep}{3.5pt}
\caption{A6 TS-Guessing contamination, by backbone.}\label{tbl:a6}
\begin{tabular}{@{}
>{\raggedright\arraybackslash}p{(\linewidth - 6\tabcolsep) * \real{0.2500}}
  >{\raggedright\arraybackslash}p{(\linewidth - 6\tabcolsep) * \real{0.2500}}
  >{\raggedright\arraybackslash}p{(\linewidth - 6\tabcolsep) * \real{0.2500}}
  >{\raggedright\arraybackslash}p{(\linewidth - 6\tabcolsep) * \real{0.2500}}@{}}
\toprule
\begin{minipage}[b]{\linewidth}\raggedright
Backbone
\end{minipage} & \begin{minipage}[b]{\linewidth}\raggedright
n diseases
\end{minipage} & \begin{minipage}[b]{\linewidth}\raggedright
Spearman \ensuremath{\rho}
\end{minipage} & \begin{minipage}[b]{\linewidth}\raggedright
Interpretation
\end{minipage} \\
\midrule
Gemini 3 Flash & 244 & \textbf{0.365} & weak positive \\
GPT-5 (reasoning=minimal) & 87 & \textbf{0.354} & weak positive \\
DeepSeek V4-Flash & 244 & \textbf{0.348} & weak positive \\
DeepSeek V4-Pro & 179 & \textbf{0.294} & weak positive \\
LIRICAL (classical Bayesian) & 26 & \ensuremath{-}0.155 & null (methodological control) \\
VC-RDAgent (offline IC) & 26 & \ensuremath{-}0.059 & null (methodological control) \\
\bottomrule
\end{tabular}
\end{table*}

\textbf{Clean dichotomy}: 4/4 LLM backbones at \ensuremath{\rho} \ensuremath{\approx} 0.29--0.37; 2/2 classical baselines
at \ensuremath{\rho} \ensuremath{\approx} 0. The classical baselines do not consume text and therefore cannot
have been ``trained on'' PubMed; their null \ensuremath{\rho} confirms our pipeline introduces
no spurious correlation. The LLM signal is real but bounded --- \ensuremath{\rho}\textsuperscript{2} \ensuremath{\approx} 0.09
means pre-cutoff exposure explains \ensuremath{\approx} 9 \% of R@1 variance, leaving \ensuremath{\approx} 91 \% to
phenotype-disease reasoning + extraction + scaffold design.

Reading is treated in \Cref{sec-7-10} (a finer-grained interpretation paired with F1);
the L4 post-cutoff PMC OA holdout (\textasciitilde200 cases, doctor-annotated in progress)
remains the bias-free reference for the residual 9 \%.

Visualised in \Cref{fig:fig4_a6_contamination_scatter}:
one panel per backbone. The clean dichotomy is visually obvious --- the
LLM backbones show a positive slope; the classical baselines
(LIRICAL, VC-RDAgent) show a flat cloud.

\subsection{Ablation A8 / H6 --- Thinking-mode (reasoning on vs off)}\label{sec-8-10}

\textbf{Question (anticipated objection: ``you crippled the models by disabling reasoning'')}:
Our main matrix runs all reasoning backbones in their minimal/off configuration
(\Cref{sec-5-2}) for cross-backbone consistency, isolation of the scaffolding effect, and
tractability. Does turning reasoning on actually help? We test this on the
single-call LLM control (no scaffold confound) with DeepSeek-V4-Pro --- the one
backbone whose reasoning can be cleanly toggled via \texttt{reasoning=\{"enabled":\ \ldots{}\}}
(GPT-5 minimal is already near-floor; V4-Pro ignores every softer throttle,
\Cref{sec-5-2} Methods note 2). Same cases, same prompt, only the reasoning flag changes.

\begin{table*}[t]
\centering\scriptsize\setlength{\tabcolsep}{3.5pt}
\caption{H6 reasoning on/off ablation (paired, PP-Store).}\label{tbl:h6}
\begin{tabular}{@{}
>{\raggedright\arraybackslash}p{(\linewidth - 6\tabcolsep) * \real{0.2500}}
  >{\raggedright\arraybackslash}p{(\linewidth - 6\tabcolsep) * \real{0.2500}}
  >{\raggedright\arraybackslash}p{(\linewidth - 6\tabcolsep) * \real{0.2500}}
  >{\raggedright\arraybackslash}p{(\linewidth - 6\tabcolsep) * \real{0.2500}}@{}}
\toprule
\begin{minipage}[b]{\linewidth}\raggedright
Config
\end{minipage} & \begin{minipage}[b]{\linewidth}\raggedright
R@1 (paired, PP-Store)
\end{minipage} & \begin{minipage}[b]{\linewidth}\raggedright
Completion
\end{minipage} & \begin{minipage}[b]{\linewidth}\raggedright
Latency/case
\end{minipage} \\
\midrule
reasoning \textbf{OFF} (main matrix) & \textbf{0.352} & 100 \% & \textasciitilde2.5 s \\
reasoning \textbf{ON} (thinking) & \textbf{0.360} & \textbf{60 \%} (40 \% no-answer) & \textasciitilde90--117 s (median), up to 571 s \\
\ensuremath{\Delta} (on \ensuremath{-} off) & \textbf{+0.008} & --- & 10--40\ensuremath{\times} slower \\
\bottomrule
\end{tabular}
\end{table*}

(N = 253 paired cases where reasoning-ON produced a parseable answer;
PP-Store; llm\_control \ensuremath{\times} V4-Pro.)

\textbf{Finding}: thinking mode changes R@1 by +0.008 --- statistically indistinguishable from zero --- while (a) failing to emit any parseable diagnosis in 40 \% of
cases (V4-Pro's unbounded reasoning consumes the entire \texttt{max\_tokens=4000} budget
before producing content) and (b) running 10--40\ensuremath{\times} slower. Reasoning-on is
therefore both \emph{unhelpful} and \emph{impractical} at benchmark scale on this task.

\textbf{Three consequences for the paper}:
1. It pre-empts the ``not tested at best'' attack: we did test thinking mode; it
does not help on retrospective rare-disease DDx from an HPO/phenotype list.
2. It justifies the reasoning-off main matrix as a design choice that loses no
accuracy while gaining tractability and cross-backbone comparability.
3. The 40 \% no-answer rate is itself a deployment-reliability finding: a
frontier reasoning model at a fixed token budget silently drops 2 in 5 cases ---
a failure mode benchmark builders and clinical integrators must budget for.

\textbf{Scope / honesty}: exploratory, single dataset (PP-Store), single control agent,
N=253 completable; the 40 \% dropout is survivorship-biased \emph{against} finding a
reasoning benefit (harder cases that need more reasoning are exactly the ones that
time out / go empty), so the true reasoning-on R@1 on all-cases could be lower, not
higher --- strengthening the ``not worth it'' conclusion rather than weakening it. We
do not extend it to the scaffolded agents because reasoning-on there is
computationally infeasible (AgentClinic reasoning-on was \textgreater900 s/case, \Cref{sec-5-2}).

